\section{Related Work}\label{sec:related}

\paragraph{Generators and relations for quantum gate groups.}
The programme of presenting quantum gate groups by generators and
relations, with completeness proved through normal forms, was opened by
Selinger's presentation of the $n$-qubit Clifford
group~\cite{selinger2015clifford}, whose staircase normal form and
``push a generator through a box'' rewrite system are the direct
ancestors of the symplectic normal form verified here.  The line
continued with presentations of $U_n(\mathbb{Z}[\tfrac12,i])$ by Bian and
Selinger~\cite{bian2021un}, building on Greylyn's
$U_4(\mathbb{Z}[1/\sqrt2,i])$~\cite{greylyn2014}; of the orthogonal
groups $O_n(\mathbb{Z}[\tfrac12])$ by Li, Ross, and
Selinger~\cite{li2021on}; of the real stabilizer operators by Makary,
Ross, and Selinger~\cite{makary2021realstabilizer}; of the $2$-qubit
Clifford+$T$ and $3$-qubit Clifford+$CS$ groups by Bian and
Selinger~\cite{bian2023cliffordt2,bian2023cliffordcs3,bian2023thesis};
of the CNOT-dihedral operators by Amy, Chen, and
Ross~\cite{amy2018cnotdihedral}; and of $3$-qubit Toffoli--Hadamard
circuits by Amy, Ross, and Wesley~\cite{amy2024toffolihadamard}.  The
first such result in an odd prime dimension is the qutrit Clifford rule
set of Li, Mosca, Ross, van de Wetering, and Zhao~\cite{li2025qutrit},
and its generalisation to all odd primes is the paper this work
verifies~\cite{bian2026qudit}.  In every one of these papers the
completeness proof is a normal-form argument checked by hand; the
Bian--Selinger papers are the exception in that their \emph{syntactic}
relation-simplification steps were checked in Agda.

\paragraph{Complete equational theories for circuits.}
A parallel programme, phrased in terms of PROPs rather than groups,
seeks complete equational theories for quantum circuits.  Cl\'ement,
Heurtel, Mansfield, Perdrix, and Valiron gave the first complete
equational theory for arbitrary quantum
circuits~\cite{clement2023complete}; it was extended to ancillae and
discarding~\cite{clement2024extensions} and made
minimal~\cite{clement2024minimal}; and Cl\'ement's complete theory for
Real-Clifford+$CH$ circuits comes with an appendix of 264 pages of
equational derivations~\cite{clement2026realcliffordch}.  Blake treats
many near-Clifford fragments uniformly as PROPs, transferring
completeness results and obtaining simpler and often minimal rule
sets~\cite{blake2026simpler}; gives a finite, dimension-uniform
presentation of qudit circuits via polycontrolled
PROPs~\cite{blake2026polycontrolled}; and proves completeness for
prime-dimensional phase-affine circuits through unique normal
forms~\cite{blake2026phaseaffine}.  All of these are pen-and-paper.
Behind circuit presentations of this kind stands Lafont's algebraic
theory of Boolean circuits~\cite{lafont2003boolean}, whose normal forms
for permutation circuits the circuit-level normal forms above
generalise.  On the graphical side, the ZX-calculus has complete
axiomatisations for the qubit stabilizer fragment, for Clifford+$T$,
and for the stabilizer fragment in odd prime
dimensions~\cite{backens2014zx,jeandel2018zx,booth2022qupitzx,poor2023travaganza};
as the paper we verify points out, completeness of a graphical calculus
for linear maps does not yield a complete calculus for unitary circuits.

\paragraph{Mechanised completeness.}
Machine-checked completeness proofs for quantum circuit rule sets are
rare.  Bian and Selinger's Agda
developments~\cite{bian2023cliffordt2,bian2023cliffordcs3} verify that
two \emph{syntactic} presentations are equivalent and never mention a
semantics.  Choudhury, Karwowski, and Sabry mechanised, in HoTT-Agda,
the presentation of the classical reversible language $\Pi$ as a free
symmetric rig groupoid~\cite{choudhury2022symmetries}, and Carette,
Heunen, Kaarsgaard, and Sabry partially mechanised their square-root
extension, including the completeness proof for the two-qubit Clifford
fragment~\cite{carette2024squareroots}; both are categorical coherence
arguments rather than presentations of concrete groups.  The library
this paper builds on~\cite{bian2026framework} was designed to close
exactly this gap; the present paper is its first application to a
research-level completeness theorem.  Verified quantum compilers and
program logics --- \textsc{Qwire}~\cite{paykin2017qwire},
\textsc{sqir}/\textsc{voqc}~\cite{hietala2021voqc},
VyZX~\cite{lehmann2025vyzx}, CoqQ~\cite{zhou2023coqq},
Qbricks~\cite{chareton2021qbricks} --- prove the \emph{soundness} of
rewrites against a matrix semantics, never the completeness of a rule
set; connecting our presented groups to such a semantics is the natural
next step (\S\ref{sec:conclusion}).

\paragraph{Group theory in proof assistants.}
Lean's mathlib defines presented groups as quotients of free
groups~\cite{mathlib2020,mathlib-doc-presented}, Isabelle's Archive of
Formal Proofs has free groups and the Nielsen--Schreier
theorem~\cite{breitner2010freegroups,kharim2023freegroups}, and the
Rocq/MathComp school~\cite{gonthier2013oddorder} develops finite group
theory over concrete permutation groups; none provides verified coset
enumeration for a \emph{given} finite relation set against a
\emph{given} concrete group.  The classical algorithms are described by
Sims~\cite{sims1994computation} and in the Handbook of Computational
Group Theory~\cite{holt2005handbook}, whose Proposition~2.55 is the
pen-and-paper form of the extension theorems we use.
